\chapter{Assessment of the September GPR-to-Poly2 phase}
\section{What improved, and what did not}
The development does not form a monotonic sequence in which each new controller is better on every metric. The initial FF retune improves ramp RMSE but worsens fuel reversal and reduces the number of accepted high-load cases. Decreasing uncertainty scheduling reduces reversal and rescues one large-step case, but gives up speed-error performance. Fixed lower gains greatly reduce ringing on a specific target, while the expanded four-input inverse still exhibits an unfavorable ringing/fine-settling tradeoff. R3 resolves much of that load-test behavior, but not a specifically GP-dependent advantage over matched PID. Covered-cycle reference tracking then reveals a different tradeoff between ramp and hold performance. Finally, Poly2 competes with GP under the same architecture, and local retuning improves its joint score at a new point while increasing peak error.

This pattern is not a failure of the research program. It identifies which claims survive controlled comparison. The strongest supported conclusion is that \emph{feedback coordinate design, independent acceleration action, and an explicit multi-objective gain selection matter at least as much as the inverse model's nominal prediction score}. The evidence also shows that an accuracy-oriented predictor can be a poor control coordinate if its local slope is weak or has the wrong sign.

\section{Internal validity and attribution}
Several design choices improve the credibility of the comparisons. The same frozen inverse is used in the setpoint and feedback branches. Loaded-controller trajectories are distinct from the nominal inverse-fit test. Same-gain FF ablations account for integral offsets rather than accidentally adding nominal fuel twice. Matched request profiles and load powers are replayed across comparators. Invalid complete runs are not scored using their accepted prefixes. S-function outputs and integral recurrences are independently reconstructed from traces. The polynomial substitution initially preserves gains before a separate local retune.

Important confounds remain. The base PI is inherited rather than retuned under the same search budget and objective as every challenger. Early FF retuning simultaneously changes the gain pair, so only the same-gain ablation isolates direct FF. R1--R3 are selected at the same target later used for their headline table. The 9250 rpm joint score uses a different reference controller from the R3 remedy score. No single fully crossed experiment covers every inverse, every architecture, and every gain strategy on the same test bank. Consequently, the narrative can identify robust local findings but cannot assign a unique numerical contribution to each design decision across the entire campaign.

\section{External validity and model mismatch}
All results use the same family of simulated plant dynamics and inherited sensor configuration. There is no hardware test, independent engine parameter set, stochastic sensor-noise campaign, fuel-actuator lag sweep, sensor-bias sweep, or execution-time qualification. A controller can be numerically accepted by this model and still be unsuitable on a real engine with unmodeled delays or constraints. Positive simulated surge margin is a monitored outcome, not an implemented surge-protection controller or a quantified safety margin under uncertainty.

The environmental inverse uses unloaded nominal data. Under load, the feedback query describes a nominal equivalent fuel coordinate rather than actual fuel demand. The posterior SD is conditional on that statistical model and does not bound the unknown-load mismatch. Training data spread across an ambient grid also do not guarantee reliable derivatives in all combinations of speed, acceleration, temperature, and pressure. Joint hull inclusion is a stronger geometric test than a rectangular bound, but convex combinations can include sparsely sampled regions and do not prove physical trajectory support or low posterior error.

The final 9250 rpm extrapolation deserves separate treatment. It is small in speed relative to the training interval, but ``small'' is not a theorem about inverse derivative accuracy. The local results are useful precisely because they are reported with that qualification; they should not be merged with the earlier in-hull cycle as if the final gains had inherited its validation.

\section{Metric validity and practical interpretation}
The paper uses RMSE, peak error, IAE, fine recovery, fuel reversal, slew, and surge margin because no single metric captures all desired behavior. A 1\% speed window is broad relative to the small disturbances in many accepted cases. Zero recovery means no excursion outside that window; it is not instantaneous correction of every disturbance. The 0.1 rpm window exposes a slow tail but is a research diagnostic rather than an established application requirement.

Fuel excess variation distinguishes reversals from monotonic fuel changes but depends on the selected window and sample interval. It can increase through high-frequency ripple or a few larger reversals, which have different actuator implications. During moving ambient conditions, some reversals may be necessary. No fuel-consumption benefit should be inferred from a lower reversal metric: total fuel mass would require $\int u\,dt$, which is a different quantity. The saved-trace sensitivity analysis in Section~\ref{sec:sensitivity} retains the final ringing ranking across four window lengths, but reverses the original-versus-retuned fine-recovery ordering at the 1 rpm band for 10\% load. Neither a lower reversal score nor a faster 0.1 rpm recovery entails better performance under every metric definition. A future actuator-oriented evaluation could add spectral energy, cycle counting, or bandwidth-weighted command variation, but these were not measured here.

The final gain choice illustrates the importance of constraints. It improves the declared joint RMSE/ringing objective while increasing peak speed error by 12.1\%. If peak error were an application constraint rather than a secondary observation, this choice might be unacceptable. Conversely, if fuel slew or fine settling dominates the application, a different point on the candidate tradeoff could be preferred. The paper reports the tested optimum under its actual objective, not a universally best controller.

\section{Statistical strength and reproducibility}
The tabulated traces are deterministic simulations. Repeating an identical run primarily checks reproducibility; it does not provide independent evidence for a confidence interval. No statistical significance claim is made from the small numerical differences between GP and Poly2 on one cycle. The held-out prediction split supports a nominal regression comparison, but time-neighboring samples from a simulator can remain correlated. The 200-point environmental test also lacks two grid conditions. Both facts limit claims of broad statistical generalization.

For that reason, the manuscript reports exact test definitions, accepted-run counts, rejected cases, domain diagnostics, and practical effect sizes rather than p-values. Full-precision summary files accompany the source. A manifest binds every imported evidence table and archived figure to its file hash. This supports a reviewer tracing a claim back to a run, while recognizing that reproducibility is distinct from physical validation.

\section{A defensible next experimental design}
The next study should hold the comparison protocol fixed before optimizing any controller. One useful design would cross four choices: conventional filtered PID, R3-GP, R3-P2, and their direct-FF variants, each with a comparable optimization budget and the same constraints. Tuning would use one set of operating points; validation would use different speeds, ambient points, load amplitudes, and model perturbations. A separate test would examine sensor lag and noise, because acceleration feedback and actual-versus-sensed ramp error are central mechanisms in the present evidence.

For model quality, the relevant checks should include not only fuel RMSE but also inverse derivative sign, derivative error against local plant perturbations, monotonicity where physically justified, and proximity to joint training support. If the operating task includes 9250 rpm, its nominal trajectories should be represented in a newly declared training/validation design. For scheduling, a matched constant-derating comparator would determine whether SD variation itself adds value beyond the mean gain reduction. For deployment, command limits, actuator dynamics, protection logic, and measured execution time would need explicit treatment.

These are proposed experiments, not completed results. They follow directly from the strongest unresolved questions in the current sequence rather than assuming that greater model complexity or a more elaborate scheduler must improve control.


\chapter{Lessons established before the October transfer tests}
\section{Summary of the research story}
The initial inverse GP provides accurate nominal fuel prediction, but a PI acting on its dynamic fuel error implicitly couples proportional, integral, and acceleration action. Adding constant setpoint FF does not independently improve fixed-reference load rejection when the integral offset is matched; early gains came from retuning. Uncertainty attenuation reduces ringing in measured cases but trades away tracking and is not a calibrated stability mechanism. Returning to fixed lower gains controls reversal, yet ambient expansion exposes weak or negative inverse slopes and slow fine recovery.

R3 separates and normalizes the static and acceleration paths, producing strong paired load-test improvements with guarded fallback. A matched PID nearly reproduces the target result, limiting claims of GP-specific benefit. Moving-reference tests then show that coverage, sensor lag, and the distinction between ramps and holds are decisive. Within the final covered cycle, R3 improves hold regulation and fuel smoothness but does not beat base PI in overall RMSE. Poly2 retains competitive control behavior despite poorer inverse prediction. Its final local retune improves RMSE and ringing at 9250 rpm, while increasing peak error and remaining slightly outside the original speed training range.

\section{Five takeaways}
\begin{enumerate}[label=\textbf{\arabic*.},leftmargin=*]
\item \textbf{Design the feedback coordinates, not just the predictor.} The virtual-fuel PI introduces coupled effective PID terms. R3's independent acceleration coefficient and slope normalization address that coupling more directly than reducing an aggregate regression error. The matched-PID diagnostic shows that the main benefit can belong to the feedback architecture and gains rather than to GPR itself.

\item \textbf{Attribute feedforward benefits with a matched-gain ablation.} At constant reference and ambient conditions, static FF is equivalent to an integral offset under the implemented law. With moving reference or ambient conditions it becomes time varying and can alter the response. In the covered cycle it improves hold accuracy but worsens overall RMSE relative to feedback-only R3. The value of FF is channel- and trajectory-dependent.

\item \textbf{Predictive standard deviation is not a gain-direction oracle.} The SD-down rule reduced fuel reversal and rescued one large-step case, but did not uniformly improve tracking or validity. The SD-up rule added little tracking benefit in the tested near-nominal uncertainty regime. Neither rule accounts for all loaded-plant mismatch, sensor uncertainty, or derivative error; both are empirical gain-scheduling hypotheses.

\item \textbf{Coverage and acceptance must be audited separately.} The broad cycle converged while spending substantial time outside the inverse and compressor-map envelopes. A final box-and-hull-covered cycle removed those flags but still showed sparse local support and a ramp/hold tradeoff. The latest 9250 rpm retune is locally useful but lies below the original speed range. Successful simulation is not a substitute for domain or robustness evidence.

\item \textbf{Use a multi-objective, matched validation before declaring a winner.} Poly2's poorer fuel-fit RMSE did not prevent competitive normalized-feedback performance. Its retuned gains lower local RMSE and ringing, confirmed at half load amplitude, but raise peak error relative to original Poly2 and retain 51\% higher edge-window peak fuel slew than base PI. Recovery ordering depends on the specified band, and the new gains have not yet transferred across the earlier grid. The appropriate conclusion is a qualified local improvement with a simpler inverse, not universal superiority of polynomial or GP control.
\end{enumerate}

